\documentclass[10pt]{article}

\usepackage[T1]{fontenc}
\usepackage[margin=0.72in]{geometry}
\usepackage{amsmath}
\usepackage{booktabs}
\usepackage{array}
\usepackage{graphicx}
\usepackage[protrusion=true,expansion=false]{microtype}
\usepackage{xcolor}
\usepackage{hyperref}
\usepackage{url}
\usepackage{caption}
\usepackage{float}
\usepackage{listings}

% A light-weight barrier marker.  Figures use top/bottom placement; avoiding
% forced page breaks keeps the report's information density consistent.
\newcommand{\FloatBarrier}{\par}
\captionsetup{font=small,labelfont=bf}

\definecolor{codegray}{RGB}{245,246,248}
\definecolor{codegreen}{RGB}{28,110,70}
\lstdefinestyle{reportcode}{
  basicstyle=\ttfamily\scriptsize,
  backgroundcolor=\color{codegray},
  commentstyle=\color{codegreen},
  frame=single,
  framesep=4pt,
  rulecolor=\color{black!18},
  columns=fullflexible,
  keepspaces=true,
  breaklines=true,
  showstringspaces=false,
  aboveskip=0.7em,
  belowskip=0.7em
}
\lstset{style=reportcode}

\graphicspath{{report/figures/}}
\hypersetup{colorlinks=true,linkcolor=blue!55!black,urlcolor=blue!55!black}
\emergencystretch=2em

\title{Collective-Aware Router Microarchitecture in gem5 Garnet\\
\large Tree multicast and congestion-aware express-link bypass}
\author{Chunyu Liu and Boyan Pu}
\date{September 6, 2026}

\begin{document}
\maketitle

\begin{abstract}
Lab~4 Topic~3 asks for two router mechanisms: bypass links that skip
intermediate Routers and multicast that sends one packet to several
destinations.  We implement both as physical Garnet datapaths rather than
traffic-generator approximations.  Tree multicast carries a destination
bitmap, replicates only at branch Routers, and reduces internal-link flits by
39.51\% on average across 162 paired cases.  Its atomic fanout nevertheless
causes ten throughput regressions, exposing a concrete scheduling tradeoff.
For bypass, a 6,144-pair source-only screen shows that hop count alone is
misleading: distance-scaled stride links are neutral at 0.999$\times$ latency
speedup.  A cycle-aware multi-hop route oracle plus conservative congestion
admission raises the same stride design to 1.430$\times$ geometric-mean
speedup and 21.87\% higher throughput over 1,536 pairs; the median is
1.082$\times$ and the worst case is 0.909$\times$.  As an additional Topic~4
case study, multi-flit tensor streaming shortens a scaled eight-H100 trace
window by 4.728$\times$ without changing source or Router flit counts.  The
report separates these claims, states their baselines and units, and retains
negative cases rather than reporting only favorable workloads.
\end{abstract}

\section{Introduction}

Modern training gives the network a repeated synchronization contract.  NCCL
exposes collectives such as broadcast and all-reduce across GPU ranks
\cite{nccl-guide}; communication therefore contains structure that ordinary
destination-by-destination unicast does not expose.  This project asks three
questions, with the first two directly matching Topic~3:

\begin{enumerate}
  \item Can Router-side replication reduce physical multicast traffic without
        changing the delivered destination set?
  \item Do express links still improve latency after their span-dependent wire
        delay, added radix, and congestion are counted?
  \item As an extension, can a traced all-reduce remain a multi-flit request
        instead of thousands of serialized scalar operations?
\end{enumerate}

Our answer is deliberately not ``all mechanisms always win.''  Multicast has
a robust traffic benefit but an atomic-fanout throughput cost in a small set
of cases.  Bypass improves loaded structured traffic, but its source-only
stride form fails once wire delay is modeled.  That failure motivates the
multi-hop and pressure-aware refinement rather than being hidden.  Tensor
all-reduce is kept as a separate Topic~4 extension so that it cannot inflate
the evidence for the two required Topic~3 mechanisms.

The implementation has three corresponding parts:

\begin{itemize}
  \item \textbf{Router-level tree multicast} shares common route prefixes and
    replicates a single destination-bitmap packet at branch Routers.
  \item \textbf{Physical express-link bypass} adds shortcut links and uses an
    offline route table to select cycle-saving express hops; the refined
    stride design may use several monotonic hops and applies congestion-aware
    admission at every eligible Router.
  \item \textbf{Tensor all-reduce replay} keeps a measured collective as a
    multi-flit request, reduces lanes in Routers, and broadcasts the result.
\end{itemize}

The evaluation follows the same structure.  Multicast is compared with
replicated unicast on the same Mesh.  Bypass is compared with plain Mesh XY
under identical packet streams and seeds.  The broad source-only screen is the
primary topology experiment; the multi-hop experiment is a like-for-like
refinement prompted by its neutral stride result.  Each claim uses its native
unit: link-flit reduction, paired latency ratio, throughput change, or trace
window.

Section~2 defines the mechanisms, baselines, and metrics; Section~3 maps them
to Garnet components and records the correctness contract.  Section~4 evaluates
multicast and express-link bypass under controlled synthetic traffic, and
Section~5 replays the measured H100 collective trace.  The final sections
record interpretation limits, authorship, and reproducibility.

\section{Mechanisms and experimental setup}

The implementation exposes two Topic~3 mechanisms and one Topic~4 extension.
The baseline and metric definitions below fix the comparison contract before
the mechanism descriptions and quantitative results.

\subsection{Experimental setup: variants and baselines}

The implementation uses gem5 v23.0.0.1 \cite{gem5-simulator} (base revision
\texttt{77fcf26d57}) and deterministic XY routing on the cycle-level Garnet
network model \cite{garnet}.  The source-only bypass screen was rerun at
project revision \texttt{135768001a97}.  The multi-hop run began at
\texttt{1e8764cde7}; its manifest hashes match the pressure-aware sources later
committed as \texttt{ea4c3c9b0f}.  Thus the recorded source snapshot, not only
the branch name at launch, identifies the evaluated binary.  A reported ratio
always compares the same
topology, traffic, packet sequence, payload, request schedule, and seed.  The
multicast comparison uses replicated unicast as its baseline: the baseline
injects one complete packet for each destination, while the proposed path
injects one destination-bitmap packet and replicates it at branch Routers.
Both variants use the same Mesh, and every selected destination must receive a
complete packet.  The bypass comparison uses plain Mesh XY as its baseline;
the proposed Mesh\_Bypass adds physical express links and an audited route
table, while injection, workload, and ordered-VNet behavior remain unchanged.

The 4$\times$4 diagrams below are illustrative mechanism examples, not the
full topology scope.  The main bypass screen covers both 4$\times$4 and
8$\times$8 meshes, the H100 replay uses a 2$\times$4 mesh, and the multicast
performance matrix uses the 4$\times$4 setup shown in
Figure~\ref{fig:multicast-tree-flow}.

Figure~\ref{fig:multicast-tree-flow} visualizes the multicast comparison before
the metric definitions below.

\begin{figure}[tb]
\centering
\begin{minipage}[t]{0.47\linewidth}
\centering
\includegraphics[width=\linewidth]{multicast_baseline_4x4.pdf}
\par\smallskip\textbf{(a) Replicated-unicast baseline}
\end{minipage}\hfill
\begin{minipage}[t]{0.47\linewidth}
\centering
\includegraphics[width=\linewidth]{multicast_tree_4x4.pdf}
\par\smallskip\textbf{(b) Proposed tree multicast}
\end{minipage}
\caption{Illustrative 4$\times$4 multicast baseline and proposed path.}
\label{fig:multicast-tree-flow}
\end{figure}

For a paired latency $L$, internal-link flit count $F$, and throughput $Q$, we
use
\[
 S_L=L_{\mathrm{base}}/L_{\mathrm{new}},\qquad
 R_F=1-F_{\mathrm{new}}/F_{\mathrm{base}},\qquad
 I_Q=Q_{\mathrm{new}}/Q_{\mathrm{base}}-1.
\]
Thus $S_L>1$ is faster, $R_F>0$ is less link traffic, and $I_Q<0$ is a
throughput regression.  One multicast request completes only after every
selected destination receives the complete packet; $Q$ counts these logical
requests per cycle, so both variants perform the same application-level work.
We use arithmetic means for additive changes and geometric means for latency
ratios.  The latter is essential for bypass because near-saturation ratios have
a long positive tail.  The released CSV files retain full distributions,
extrema, and regression cases.

Table~\ref{tab:methodology} collects the parameters that otherwise tend to be
hidden in command lines.  Each comparison is paired by workload and seed.

\begin{table}[tb]
\centering
\caption{Main controlled-study configurations.}
\label{tab:methodology}
\small
\begin{tabular}{@{}p{0.20\linewidth}p{0.34\linewidth}p{0.34\linewidth}@{}}
\toprule
Parameter & Multicast matrix & Bypass screen \\
\midrule
Topology & 4$\times$4 Mesh, source Router 5 & 4$\times$4 and 8$\times$8 Mesh \\
Traffic & 4/8/16 destinations; optional uniform background & uniform, transpose, bit complement, hotspot \\
Packet size & 1/4/16 flits & 1/4/16/64 flits \\
Load points & 0.1/0.5/1.0, background 0/0.1 & 16 points from 0.01 to 3.0 \\
Seeds & 1, 7, 17 & 1, 7, 17 \\
Measurement & 20 warm-up, 100 measured, 20 cooldown requests & 1,000 warm-up, 5,000 measurement, 100,000 drain cycles \\
\bottomrule
\end{tabular}
\end{table}

\subsection{Tree multicast}

The baseline Network Interface injects one physical packet per destination.
The multicast path injects one packet carrying a 64-bit destination bitmap.
Each Router removes destinations already served, computes the required output
branches, and makes a local copy only where selected.  Head, body, and tail
state is retained per branch.  Fanout is atomic: a flit advances only when all
selected output VCs have credit.  This preserves exact multi-flit delivery,
but it also couples a fast branch to a blocked branch---an important
interpretation of the throughput results below.  The bitmap bounds one
collective domain to 64 Routers.

\subsection{Offline route-table express-link bypass}

\texttt{Mesh\_Bypass} appends bidirectional diagonal or stride links to the
ordinary Mesh.  Before simulation, a route-table builder enumerates every
current-Router/destination pair.  The original screen permits one source
express hop followed by XY.  The refinement lets dimension-aligned stride
links be selected repeatedly in the X phase and then the Y phase, but only
when the complete suffix has lower modeled cycle cost.  The builder constructs
the full paths and channel-dependency graph and rejects any placement whose
graph is cyclic.

\begin{itemize}
  \item \textbf{Optimistic:} every express link costs one cycle, an upper
        bound on the value of shortening a path.
  \item \textbf{Distance-scaled:} latency and serialization scale with the
        Manhattan span, a more conservative proxy for a physical wire.
\end{itemize}

The cost of an ordinary edge is link latency plus Router latency.  An express
edge uses its configured wire latency plus one Router traversal and is admitted
statically only if it lowers the full suffix cost.  Ordered Vnets retain this
deterministic table.  On unordered traffic, the refined selector may fall back
to the monotonic XY edge when the express output or landing output is blocked,
the express output is materially more congested, or the destination is detected
as a sustained hotspot.

\begin{enumerate}
  \item Each Router consults the table entry indexed by its own ID and the
        destination.
  \item A stride express entry must move monotonically toward the destination
        in the current XY dimension and may not overshoot.
  \item Both adaptive choices therefore remain in X-before-Y order.
  \item A full all-pairs channel-dependency DAG certifies the resulting order.
\end{enumerate}

The bypass screen uses data VNet 2 (four buffers per VC).  The express-link cost is
either optimistic (one cycle) or distance-scaled by the link's Manhattan span;
the latter is the hardware-relevant model used for the main comparison.

The bypass comparison therefore has a plain Mesh XY baseline and a
Mesh\_Bypass variant with physical shortcut candidates.  Figure~\ref{fig:bypass-topology-route-table}
shows the two candidate placements.  The diagonal design retains the original
source-only policy; the stride refinement can select multiple aligned express
hops along the deterministic dimension-ordered path.

\begin{figure}[tb]
\centering
\begin{minipage}[t]{0.47\linewidth}
\centering
\includegraphics[width=\linewidth]{bypass_diagonal_4x4.pdf}
\par\smallskip\textbf{(a) Diagonal placement}
\end{minipage}\hfill
\begin{minipage}[t]{0.47\linewidth}
\centering
\includegraphics[width=\linewidth]{bypass_stride_4x4.pdf}
\par\smallskip\textbf{(b) Stride-2 placement}
\end{minipage}
\caption{Illustrative 4$\times$4 bypass placements and route-table-selected
paths.}
\caption*{\textbf{Legend:} orange dashed lines are installed bidirectional
bypass links; solid orange is the selected express hop or hops; solid blue is
the ordinary DOR segment; gray lines are ordinary Mesh links.}
\label{fig:bypass-topology-route-table}
\end{figure}

\FloatBarrier
\section{Implementation and validation}

The implementation changes the packet metadata, Network Interface (NI),
Router fast path, routing unit, and topology construction rather than emulating
the mechanisms in the traffic generator.  Table~\ref{tab:implementation-map}
summarizes the division of responsibility.

\begin{table}[tb]
\centering
\caption{Implementation map in gem5 Garnet.}
\label{tab:implementation-map}
\small
\begin{tabular}{@{}p{0.16\linewidth}p{0.29\linewidth}p{0.45\linewidth}@{}}
\toprule
Mechanism & Main components & Added state and behavior \\
\midrule
Multicast & \texttt{flit}, NI, \texttt{Router} & 64-bit destination mask,
per-packet branch VCs, Router-side replication, exact completion tracking \\
Bypass & \texttt{Mesh\_Bypass.py}, \texttt{bypass\_oracle.py},
\texttt{RoutingUnit} & Physical express ports, all-pairs per-Router next-hop
table, monotonic DOR paths, channel-dependency audit, runtime fallback \\
Tensor & NI, \texttt{Router}, replay compiler & Per-flit lane ID,
per-$(\text{collective},\text{lane})$ accumulator, queued backpressure-safe
reduce/broadcast forwarding \\
\bottomrule
\end{tabular}
\end{table}

\subsection{Atomic multicast fanout}

The head flit first computes the pruned-tree children selected by its bitmap.
The Router commits branch state only if every required output has a free VC;
body and tail flits reuse those VCs.  Listing~\ref{lst:multicast} is an excerpt
from the implemented two-phase allocation.  This makes a multi-flit packet
exact and prevents partial branch allocation, but also explains the measured
head-of-line coupling: one blocked branch stalls the whole fanout.

\noindent\begin{minipage}{\linewidth}
\begin{lstlisting}[language=C++,caption={Source excerpt: atomic multicast branch allocation (\texttt{Router.cc}).},label={lst:multicast}]
const int vnet = head_flit->get_vnet();
for (const int destination : destinations) {
    const int outport = collectiveOutportForChild(destination);
    if (outport < 0 || !m_output_unit[outport]->has_free_vc(vnet)) {
        m_network_ptr->recordMulticastCreditStall();
        return false;
    }
}
for (const int destination : destinations) {
    const int outport = collectiveOutportForChild(destination);
    const int outvc = m_output_unit[outport]->select_free_vc(vnet);
    fatal_if(outvc < 0, "Router %d lost multicast VC allocation", m_id);
    m_multicast_branches.push_back({outport, outvc, destination});
}
\end{lstlisting}
\end{minipage}

\subsection{Auditable bypass routing}

Topology construction creates two directed Garnet links for every installed
physical shortcut and scales their latency by Manhattan span in the
distance-scaled model.  Listing~\ref{lst:bypass} shows the implemented dynamic
programming recurrence: it compares complete ordinary and express suffixes,
and accepts an express edge only when \texttt{dor\_express} proves monotonic
X-before-Y progress.  Exact-cost ties prefer the ordinary edge, so a bypass
must save a modeled cycle rather than merely a graph hop.

\begin{lstlisting}[language=Python,caption={Source excerpt: cycle-aware multi-hop DOR selection (\texttt{bypass\_oracle.py}).},label={lst:bypass}]
@lru_cache(maxsize=None)
def best_stride_step(current, destination):
    if current == destination:
        return 0, None
    ordinary_next = xy_step(current, destination)
    suffix_cost, _ = best_stride_step(ordinary_next, destination)
    choices = [(ordinary_edge_cost + suffix_cost, 0,
                ordinary_next, "", None)]
    for link in outgoing[current]:
        if not dor_express(link, destination):
            continue
        suffix_cost, _ = best_stride_step(link["dst"], destination)
        cost = int(link["latency"]) + router_latency + suffix_cost
        choices.append((cost, 1, link["dst"], link["name"], link))
    cost, _, _, _, selected = min(choices)
    return cost, selected
\end{lstlisting}

The builder then materializes every source--destination path, adds each
successive channel pair to a dependency graph, and requires a topological
ordering over the entire graph.  This is stronger than checking that each
individual path terminates: a cycle formed across several routes also fails
construction.  Runtime adaptation can choose only the table's express edge or
the monotonic XY edge in the same DOR phase, so it does not introduce a new
channel class or a dimension reversal.

The performance refinement is completed in \texttt{RoutingUnit.cc}, not in
the offline oracle alone.  After computing local VC/credit pressure and the
landing Router's oracle-selected next output, the retained policy executes the
fallback in Listing~\ref{lst:bypass-admission}.  A sustained many-to-one
destination, a blocked express/escape output, or a materially healthier XY
output rejects the shortcut.  Long packets additionally require actual credit
headroom on both the express and landing outputs.

\begin{lstlisting}[language=C++,caption={Source excerpt: conservative runtime bypass admission (\texttt{RoutingUnit.cc}).},label={lst:bypass-admission}]
if (network->bypassDestinationSkewed(route.dest_router) ||
    express_free == 0 || escape_free == 0 ||
    (long_packet && lookahead_free == 0) ||
    express_free + 1 < xy_free)
    return xy_outport;

const int landing_capacity = vcs_per_vnet *
    network->getBuffersPerDataVC();
if (long_packet &&
    ((xy_free == vcs_per_vnet && xy_credits == landing_capacity) ||
     express_free < xy_free || express_credits < xy_credits ||
     lookahead_free < vcs_per_vnet ||
     lookahead_credits < landing_capacity))
    return xy_outport;
\end{lstlisting}

\subsection{Streaming tensor reduction}

Tensor packets retain their multi-flit boundary while each flit carries a lane
identifier.  Listing~\ref{lst:tensor} is the Router accumulation point.  An
entry is released as soon as every local/child contribution arrives; completed
lanes wait in a FIFO when their next output VC is unavailable, so backpressure
does not discard reduction state.

\begin{lstlisting}[language=C++,caption={Abridged source excerpt: lane-wise tensor accumulation (\texttt{Router.cc}).},label={lst:tensor}]
const int lane = t_flit->get_lane_id();
const auto key = std::make_pair(t_flit->get_collective_id(), lane);
CollectiveLaneState& state = m_collective_lanes[key];
state.sum += t_flit->get_value();
++state.count;
if (state.count == m_collective_expected_fanin) {
    const int64_t sum = state.sum;
    m_collective_lanes.erase(key);
    if (m_collective_root)
        forwardCollectiveLane(sum, lane, CollectiveOp::Broadcast,
                              m_id, t_flit);
    else
        // parent is derived from m_collective_parent_outport
        forwardCollectiveLane(sum, lane, CollectiveOp::Reduce,
                              parent, t_flit);
}
\end{lstlisting}

\subsection{Validation contract}

Correctness gates check payload values and delivery sets, not only simulator
termination.  The compact-data builder additionally rejects incomplete pairs,
non-finite metrics, inconsistent route metadata, and wrong run counts before
generating any headline figure.

\begin{table}[tb]
\centering
\caption{Executed validation coverage in the accepted snapshot.}
\label{tab:validation}
\small
\begin{tabular}{@{}lrl@{}}
\toprule
Gate & Coverage & Required invariant \\
\midrule
Multicast matrix & 208 paired cases & Exact bitmap, flit order, no missing/duplicate delivery \\
Bypass screen/refinement & 10,752 runs & Complete table, route termination, acyclic CDG \\
Multicast $\times$ bypass & 624 runs & 416 paired interactions complete without route corruption \\
Tensor functional & 14 cases & Exact lane count and reduced value on varied meshes \\
Tensor backpressure & 32 cases & Completion under varied link/router latency and occupancy \\
Trace tools & 12 unit tests & Schema, sizes, ordering, IDs, and participant validation \\
\bottomrule
\end{tabular}
\end{table}

\section{Evaluation results}

\subsection{Multicast performance}

Tree multicast removes physical link work reliably, but its throughput benefit
depends on branch synchronization.  The results below separate those two
claims rather than treating traffic reduction as a proxy for throughput.

Across 162 paired cases, tree multicast reduces internal-link flits by
39.51\% on average (median 41.18\%, range 16.67--53.13\%).  The effect is
monotone in fanout: 26.19\% for four destinations, 39.22\% for eight, and
53.13\% for sixteen.  This is the cleanest evidence for the mechanism because
both variants use the same Mesh and the metric counts only physical link
flits.  Figure~\ref{fig:multicast-traffic} summarizes the reduction by
fanout.

\begin{figure}[tb]
\centering
\includegraphics[width=0.84\linewidth]{traffic_reduction_by_group.pdf}
\caption{Internal-link flit reduction from tree replication; bars are means and
whiskers show the observed range for each fanout.}
\label{fig:multicast-traffic}
\end{figure}

Latency also improves in this matrix: geometric-mean speedup is 3.638$\times$ (median
3.588$\times$, range 1.200--22.512$\times$), and all 162 paired values exceed
one.  The multicast implementation uses a dedicated Router fast path, so its
switch-allocation and crossbar sequence differs from unicast.  We therefore
read latency as an implemented-system result and use link flits for the
structural comparison.

Lower link traffic does not always translate into higher throughput.  The
arithmetic mean change is
+190.89\%, but the distribution is highly skewed by full-group traffic; ten
of 162 cases regress, all at four destinations, with a worst case of
$-18.24\%$.  The paired distribution is shown in
Figure~\ref{fig:multicast-throughput}.  Atomic fanout explains how both
statements can be true: the tree does less total link work, yet a single
blocked branch can hold up all copies.

\begin{figure}[tb]
\centering
\includegraphics[width=0.84\linewidth]{throughput_improvement_distribution.pdf}
\caption{Paired multicast throughput changes, including the high-fanout tail
and the ten regression cases.}
\label{fig:multicast-throughput}
\end{figure}

\subsection{Multicast sensitivity and tradeoff}

The workload breakdown makes the dependency concrete.  At 4, 8, and 16
destinations, mean throughput changes are +34.38\%, +154.87\%, and +383.42\%,
but the four-destination group contains all ten regressions.  By packet size,
the mean changes are +261.65\% (1 flit), +169.45\% (4 flits), and +141.57\%
(16 flits); the negative cases concentrate in the longer, loaded four-way
requests.  Figure~\ref{fig:multicast-scaling} shows the same pattern together
with the min--max spread, while Figure~\ref{fig:multicast-tradeoff} relates
traffic reduction to throughput on a case-by-case basis.

\begin{figure}[tb]
\centering
\includegraphics[width=0.88\linewidth]{multicast_scaling.pdf}
\caption{Sensitivity of multicast throughput and latency to fanout, packet
size, and offered load.}
\label{fig:multicast-scaling}
\end{figure}

\begin{figure}[tb]
\centering
\includegraphics[width=0.82\linewidth]{multicast_traffic_throughput_tradeoff.pdf}
\caption{Traffic reduction versus throughput change for individual paired
cases; the dashed line marks unchanged throughput.}
\label{fig:multicast-tradeoff}
\end{figure}

\paragraph{Multicast takeaway.}
Tree replication reliably removes duplicate physical work, and the saving
grows with fanout.  Throughput follows a different workload curve because
atomic fanout and the dedicated fast path introduce their own scheduling
effects.

\FloatBarrier
\subsection{Primary experiment: offline route-table bypass}

Express-link bypass is useful only when path shortening survives wire cost and
contention.  The screen below tests that claim across topology, traffic, and
wire-model choices.

The bypass sweep uses four standard synthetic traffic patterns.  In
\texttt{uniform\_random}, each source chooses a destination uniformly at
random.  In \texttt{transpose}, a source at mesh coordinate $(x,y)$ sends to
$(y,x)$, creating a structured permutation.  In \texttt{bit\_complement}, it
sends to the mirrored coordinate $(W-1-x,H-1-y)$, pairing opposite locations
in the mesh.  In \texttt{hotspot}, 50\% of packets target one fixed Router
(Router $N/2$ in an $N$-node mesh) and the rest use uniform random
destinations.  These four patterns isolate route structure and contention.

The original offline route-table sweep compares the same synthetic workloads
and seeds on Mesh XY and on each physical express-link design, with the route
table selecting at most one source express hop.  The latency column in
Table~\ref{tab:bypass-screen} is the geometric mean over all 1,536 paired seed
cases per design; additive changes use arithmetic means.  Reporting the
median and regression tail later is necessary because saturated baseline
points can produce very large ratios.  The fixed table isolates the value of
an admissible shortcut, while the regression count exposes how often the
aggregate result fails to transfer to an individual workload.

\begin{table}[ht]
\centering
\caption{Original source-only topology and wire-model screen over 1,536
paired seed cases per design; latency ratios use geometric means.}
\label{tab:bypass-screen}
\begin{tabular}{@{}lrrrr@{}}
\toprule
Design & Latency speedup & Regressions & Throughput change & Router-work change \\
\midrule
Diagonal, distance-scaled & 1.131$\times$ & 350/1,536 & +9.41\% & +5.49\% \\
Diagonal, optimistic & 1.148$\times$ & 164/1,536 & +9.44\% & +5.53\% \\
Stride, distance-scaled & 0.999$\times$ & 539/1,536 & +2.72\% & $-$11.35\% \\
Stride, optimistic & 1.024$\times$ & 321/1,536 & +2.81\% & $-$11.26\% \\
\bottomrule
\end{tabular}
\par\medskip\footnotesize Router-work change is
$F_{\mathrm{new}}/F_{\mathrm{base}}-1$ for each run's \emph{aggregate} Router
traversals.  Positive throughput can therefore produce a positive traversal
count even when individual paths shorten; negative values mean less Router
work.  The throughput and Router-work columns must therefore be read together.
\end{table}

The cost records explain the design choice.  A 4$\times$4 diagonal placement
adds 9 undirected links (wire proxy 18, maximum radix 8), while stride adds 16
(wire proxy 32, maximum radix 6); on 8$\times$8 the corresponding link counts
are 49 and 96.  Physical topology choice therefore combines hop reduction with
wire length, radix, and buffer cost.

Figure~\ref{fig:bypass-overall} adds the multi-hop refinement to the four
source-only designs.

\begin{figure}[tb]
\centering
\includegraphics[width=0.88\linewidth]{bypass_overall.pdf}
\caption{Latency speedup and throughput change for the source-only screen and
the distance-scaled multi-hop adaptive refinement.}
\label{fig:bypass-overall}
\end{figure}

In the source-only screen, diagonal links provide the strongest latency result,
while stride links remove more Router traversals but incur a larger
physical-link budget.  At
offered load $\leq 0.1$, however, the distance-scaled geometric means are
0.999$\times$ for diagonal and 0.993$\times$ for stride.  The gain therefore
comes from favorable loaded cases rather than a universal reduction in base
latency.  The optimistic wire model is a sensitivity bound; distance-scaled
latency is the hardware-relevant comparison.

\subsection{Multi-hop stride refinement}

The neutral source-only stride result identifies a structural limit: a packet
can skip only one Router even when its path crosses several aligned shortcuts.
The refined policy consults the cycle-aware table at every Router, so a packet
may use several stride links while remaining monotonic in X and then Y.
For unordered data traffic, conservative admission retains the static shortcut
unless the express or landing output is blocked, the express output is
materially more congested than XY, or epoch counters identify the destination
as a sustained hotspot.  For packets above 32 flits, the selector also checks
credit headroom and the oracle-selected output at the landing Router, admitting
the shortcut only when XY is under pressure and the landing output is idle.
A 96-pair pilot rejected a more aggressive version despite its higher mean
because it introduced four regressions larger than 5\%; the retained guard had
none in the same pilot.

\begin{table}[tb]
\centering
\caption{Distance-scaled multi-hop stride refinement.  Latency is a geometric
mean; throughput and Router-traversal-per-delivered-flit changes are arithmetic
means.  The last column counts regressions larger than 5\%.}
\label{tab:bypass-multihop}
\small
\begin{tabular}{@{}lrrrrr@{}}
\toprule
Group & Pairs & Speedup & Throughput & Router/flit reduction & $>5\%$ regress. \\
\midrule
Overall & 1,536 & 1.430$\times$ & +21.87\% & 18.10\% & 1 \\
4$\times$4 & 768 & 1.291$\times$ & +4.37\% & 17.37\% & 1 \\
8$\times$8 & 768 & 1.583$\times$ & +39.37\% & 18.84\% & 0 \\
Uniform random & 384 & 1.842$\times$ & +19.63\% & 24.98\% & 0 \\
Transpose & 384 & 1.182$\times$ & +5.50\% & 15.92\% & 0 \\
Bit complement & 384 & 1.892$\times$ & +61.07\% & 14.54\% & 0 \\
Hotspot & 384 & 1.015$\times$ & +1.28\% & 16.99\% & 1 \\
\bottomrule
\end{tabular}
\end{table}

Across all 1,536 pairs, the median is 1.082$\times$ and 243 ratios are below
one, but only one regresses by more than 5\%.  The best ratio, 97.94$\times$,
comes from a near-saturated baseline and is why we do not report an arithmetic
mean as the headline.  The worst is 0.909$\times$ for
4$\times$4, 16-flit hotspot traffic at offered load 0.3.  At offered load at
most 0.1 the geometric mean is 1.051$\times$; above 0.7 it is
1.627$\times$.  The 64-flit group now reaches 1.134$\times$ with no regression
larger than 5\%.  Thus the 1.430$\times$
headline is not a claim of universal base-latency improvement: multi-hop
bypass is most valuable when the skipped Router stages also relieve loaded
paths, while hotspot admission largely restores the baseline on
many-to-one traffic.

We also tested a natural follow-up rather than silently tuning around that
worst case.  On the same 96-pair pilot, disabling bypass only within Manhattan
radius 1, 2, or 3 of a detected hotspot produced overall geometric means of
1.6830$\times$, 1.6838$\times$, and 1.6837$\times$, all below the retained
global guard's 1.6844$\times$.  More importantly, the targeted worst case was
0.9001$\times$, 0.9090$\times$, and 0.9090$\times$, below the predeclared
0.95$\times$ acceptance threshold.  We therefore rejected all three before a
full run.  This negative pilot is retained in
\path{report/data/bypass_hotspot_radius_pilot.csv} and is not mixed into the
1,536-pair headline result.

\FloatBarrier

\paragraph{Bypass takeaway.}
Express links are useful when route structure makes a saved hop valuable, while
distance-scaled wires and extra radix impose real costs.  Moving from one
source-only shortcut to audited multi-hop DOR raises geometric-mean speedup
from 0.999$\times$ to 1.430$\times$ for stride links; congestion admission
reduces the severe-regression tail without weakening the wire model.

\FloatBarrier
\subsection{Multicast--bypass interaction}

A 624-run factorial gate combines both multicast modes with Mesh XY and the
bypass topology.  All 416 Mesh/bypass pairs complete with exact destination
sets and unchanged wire-distance accounting.  Table~\ref{tab:interaction}
shows that shortcut use remains mechanism-dependent: replicated unicast can
route each copy independently, whereas tree multicast preserves its pruned XY
branch semantics.  This gate establishes composability; it is not substituted
for the broader unicast bypass screen above.

\begin{table}[tb]
\centering
\caption{Multicast--bypass interaction over 104 pairs per row.}
\label{tab:interaction}
\small
\begin{tabular}{@{}llrr@{}}
\toprule
Multicast mode & Wire model & Link-flit reduction & Latency regressions \\
\midrule
Replicated unicast & Distance-scaled & 16.10\% & 39/104 \\
Replicated unicast & Optimistic & 16.10\% & 0/104 \\
Tree multicast & Distance-scaled & 8.33\% & 15/104 \\
Tree multicast & Optimistic & 8.33\% & 0/104 \\
\bottomrule
\end{tabular}
\end{table}

\section{H100 trace case study}

This additional Topic~4 case study preserves a measured collective's ordering
while testing request representation.  Its scaled tick counts remain separate
from the factorial Topic~3 aggregates.

\subsection{Trace lowering and implementation}

The input is an executed eight-H100 80GB HBM3 NCCL microbenchmark
(PyTorch 2.8.0+cu128, CUDA 12.8, NCCL 2.27.3).  The replay
compiler selects the 15 measurement-phase broadcasts and 15 all-reduces,
represents communication with 16-byte flits, and scales payload bytes and
relative release times by $1/1024$.  This preserves event ordering while
making the trace practical for cycle-level replay.

For broadcast, each selected source event is lowered to variable-length
tree-multicast packets.  For all-reduce, the scalar replay creates independent
serialized lanes; each lane converges rank values in the Routers and then
travels back down the convergence tree.  The tensor replay keeps each event as
one multi-flit logical request, tracks a lane ID per flit, and performs the
same Router-side reduction and tree broadcast while retaining the request's
multi-flit boundary.  The scalar and tensor paths are distinct collective
implementations from the bitmap-based \texttt{tree\_multicast} path; all three
exploit shared tree prefixes.

Figure~\ref{fig:tensor-allreduce-pipeline} connects the datapath to the replay
representation.  It also makes clear why the tensor comparison preserves the
same contribution and Router-flit counts: only the logical request boundaries
and completion scheduling change.

\begin{figure}[tb]
\centering
\includegraphics[width=0.90\linewidth]{tensor_allreduce_pipeline.pdf}
\caption{Tensor all-reduce datapath and the scalar-lane and tensor-request
replay lowerings.}
\label{fig:tensor-allreduce-pipeline}
\end{figure}

The replay compares Mesh XY with a four-link, distance-scaled diagonal bypass
arm.  Every route selected in this 2$\times$4 placement is ordinary XY, making
the bypass arm a noninterference control rather than evidence of shortcut
speedup.

\subsection{Replay experiments and results}

The broadcast replay expands 15 source events into 30 tree-multicast requests;
both designs complete 6,720 source flits and 210 destination deliveries in a
17,047,500-tick window.  The all-reduce replay compares serialized scalar lanes
with multi-flit tensor requests under the same route and release schedule.

Table~\ref{tab:h100-replay} records the three replay views.  Both topology arms
are identical, confirming that unused physical links do not perturb the trace.

\begin{table}[ht]
\centering
\caption{Scaled H100 replay summary; p95 values use different logical units
and must not be compared across scalar and tensor rows.}
\label{tab:h100-replay}
\begin{tabular}{@{}lrrrr@{}}
\toprule
Replay view & Logical requests & Source flits & Window (ticks) & p95 (ticks) \\
\midrule
Broadcast & 30 & 6,720 & 17,047,500 & 643,500 \\
Scalar all-reduce lanes & 6,720 & 53,760 & 80,638,000 & 10,000 \\
Tensor all-reduce & 15 & 53,760 & 17,054,500 & 2,567,500 \\
\bottomrule
\end{tabular}
\end{table}

Both physical designs complete 53,760 exact lane reductions and deliveries
with 147,840 Router flits.  Tensor-request p50/p95/p99 latency is
647,500/2,567,500/2,567,500 ticks; these values cannot be compared directly
with scalar-lane p95 because the logical request units differ.

\begin{figure}[tb]
\centering
\includegraphics[width=0.74\linewidth]{tensor_measurement_comparison.pdf}
\caption{End-to-end replay window for scalar-lane and tensor-request
lowerings of the same scaled all-reduce trace.}
\label{fig:tensor}
\end{figure}

Scalar lowering records an 80,638,000-tick window, so tensor streaming is
4.728$\times$ shorter on both Mesh XY and Mesh Bypass.  Source and Router
flit counts are identical, which isolates the measured gain to request
representation and completion scheduling while rank contributions and network
traffic remain fixed.  The result characterizes the configured replay model;
translating it to training-step speed requires an end-to-end study.

\paragraph{Trace-replay takeaway.}
Keeping a tensor as a streaming collective removes scalar serialization in
this network model while preserving exact arithmetic and delivery invariants.
The effect appears in the total trace window; per-request tail percentiles use
different logical units for the two lowerings.

\FloatBarrier

\section{Limitations and interpretation boundary}

The experiments establish the mechanisms within the tested model.  The scope
of the evidence is:

\begin{itemize}
  \item The multicast matrix covers one 4$\times$4 source placement (Router
        5); corner, edge, and center source placements are future coverage.
        Its dedicated fast path and atomic fanout expose branch serialization
        under load.
  \item Bypass wire cost uses a Garnet latency, serialization, and static-cost
        proxy.  Area, power, repeaters, and timing closure remain outside the
        model.  Refined stride routes may use multiple monotonic express hops;
        their complete channel-dependency graph is audited before simulation.
  \item The bypass screen uses synthetic unicast traffic on data VNet 2.  The
        distance-scaled screen is the hardware-oriented comparison, with
        optimistic timing serving as a sensitivity point.
  \item The H100 input is an executed collective microbenchmark with
        $1/1024$ byte and time scaling.  Full model traces are future workload
        coverage.  Broadcast and tensor replay on the shared topology use
        ordinary XY paths.
  \item Tensor accumulation uses software state with arithmetic latency and
        finite table capacity outside the model.  Chunk boundaries are
        validated by the contract and carried by the replay protocol.
  \item The results characterize cycle-level Garnet mechanisms.  NVLink/NVSwitch
        measurements and end-to-end model throughput are future validation
        targets.
\end{itemize}

\section{Division of labor}

Chunyu Liu authored multicast and bypass implementation, correctness checks, trace capture,
optimization, and the associated analysis.  Boyan Pu authored the tensor
all-reduce datapath, replay contract, validator, backpressure stress cases, and
tensor statistics.  Both authors integrated branches, reviewed regressions,
reran the final quantitative snapshot, and verified this report.

\section{Reproducibility and artifact trail}

Compact accepted snapshots are in \path{report/data/}; generated vector and
preview figures are in \path{report/figures/}; raw outputs are git-ignored
under \path{report/raw-results/head-77fcf26d/}.  The source trace and replay
hashes, software versions, scaling semantics, and artifact roots are recorded
in \path{report/data/trace_replay_provenance.txt}.  The compact-data builder
rejects wrong run counts, non-finite values, failed regression stages, inconsistent VNet or
routing metadata, incomplete multicast pairs, and missing designs before
writing summary tables or figures.

The original source-only bypass matrix contains 7,680 simulator runs.  The
accepted multi-hop refinement adds 3,072 runs and stores its exact arguments,
simulation windows, and source hashes in
\path{report/data/bypass_multihop_manifest.json}.  The resulting 1,536 paired
rows are summarized in \path{report/data/bypass_multihop_summary.csv}; the
policy-selection and rejected-radius pilots are retained separately, so they
cannot be mistaken for final-matrix samples.

The full command matrix is retained in \path{report/README.md}.  The compact
snapshot and every figure can be rebuilt with:

\begin{verbatim}
python3 tests/gem5/lab4/run_lab4_full_gate.py --jobs 64
python3 report/scripts/build_final_report_data.py
python3 report/scripts/plot_multicast.py
python3 report/scripts/generate_architecture_figures.py
\end{verbatim}

The trace-replay runners retain their original four-link diagonal topology as
a controlled H100 case-study arm.  The new multi-hop stride result belongs to
the broad synthetic screen and is not retroactively substituted into that
fixed trace placement.

\section{Conclusion}

Both required Topic~3 mechanisms are implemented and evaluated against matched
baselines.  Tree multicast shares route prefixes and removes 39.51\% of
internal-link flits on average, rising to 53.13\% at fanout 16.  Its ten
throughput regressions show that atomic branch synchronization remains a real
microarchitectural cost.  For bypass, the source-only distance-scaled screen
is the diagnostic result: diagonal reaches 1.131$\times$, while stride is
neutral at 0.999$\times$ despite reducing Router traversals.  Reconsulting an
audited cycle-cost oracle at every Router and admitting express hops only under
safe local and landing pressure raises stride to 1.430$\times$ geometric-mean
speedup and +21.87\% throughput.  The 1.082$\times$ median, 0.909$\times$
worst case, and rejected radius sweep bound that claim.  The additional tensor
experiment preserves 15 multi-flit requests and shortens the scaled H100
window by 4.728$\times$ with identical network work.

The common lesson is narrower than ``more network features are faster.''
Collective structure should be exposed only where it removes measurable
physical work; a shortcut should be selected by modeled cycle cost and current
pressure rather than hop count; and every speedup should be paired with its
distribution, regression tail, and completion contract.

\begin{thebibliography}{9}

\bibitem{gem5-simulator}
N. Binkert et al., ``The gem5 Simulator,'' \emph{ACM SIGARCH Computer
Architecture News}, vol. 39, no. 2, 2011, pp. 1--7.
doi: \href{https://doi.org/10.1145/2024716.2024718}{10.1145/2024716.2024718}.

\bibitem{garnet}
N. Agarwal, T. Krishna, L.-S. Peh, and N. K. Jha, ``GARNET: A Detailed
On-Chip Network Model inside a Full-System Simulator,'' \emph{ISPASS}, 2009,
pp. 33--42.
doi: \href{https://doi.org/10.1109/ISPASS.2009.4919636}{10.1109/ISPASS.2009.4919636}.

\bibitem{nccl-guide}
NVIDIA, ``NVIDIA Collective Communication Library Documentation.''
\url{https://docs.nvidia.com/deeplearning/nccl/user-guide/}.

\end{thebibliography}

\end{document}
